\section{Introducción: la AI Psychosis y el cambio radical de paradigma en la programación}

Karpathy usa una palabra muy gráfica para describir el estado actual: \textbf{AI Psychosis} (psicosis de la IA). Desde diciembre de 2024, pasó de una proporción de 80\%/20\% entre código escrito a mano y código delegado a agentes, a que casi el 100\% lo haga un agente. Incluso dice que «desde diciembre, es posible que no haya escrito ni una sola línea de código con mis propias manos».

\begin{importantbox}{El momento clave del cambio de paradigma}
En diciembre de 2024, la capacidad de los coding agents dio un salto. Karpathy lo describe como «something flipped» (algo se invirtió): el flujo de trabajo diario de los ingenieros de software cambió por completo de la noche a la mañana. Es posible que la mayoría todavía no se haya dado cuenta de que esto ocurrió, pero para quienes desarrollan en la frontera, se trata de un cambio de paradigma irreversible.
\end{importantbox}

\begin{knowledgebox}{De la velocidad de escritura al rendimiento de tokens}
Antes, el cuello de botella de la productividad individual era la velocidad de escritura y la habilidad personal para programar. Ahora, el cuello de botella se ha vuelto «cuánto rendimiento de tokens puedes despachar». Karpathy compara esto con la ansiedad de la época del doctorado en torno a la utilización de las GPU: «la ansiedad de que las GPU estén ociosas» se convirtió en «la ansiedad de no agotar la cuota de tokens».
\end{knowledgebox}

